% Solutions/hints for ch25 exercises (assembled into Appendix C)

\begin{solution}{exr:ch25-metrics-log}
(a) Drone--drone: the sampled distances are $2$, $\sqrt5$, $\sqrt5$, $\sqrt8$ and the segment minima are no smaller, so $d_{\min}^{\mathrm{dd}} = 2.0$~m. Drone~1 against the intruder: samples $\sqrt{20}, \sqrt{10}, 2, \sqrt2$; on the last interval the relative position goes from $(0,-2)$ to $(1,-1)$, $s^* = 1$, so the segment minimum is again $\sqrt2 = 1.414$~m. Drone~2 against the intruder: samples $\sqrt8, \sqrt5, 1, \sqrt2$, and on the interval from $t=2$ to $t=3$ the relative position goes from $(-1,0)$ to $(-1,1)$ with $s^* = 0$: minimum $1.0$~m. Hence $d_{\min}^{\mathrm{di}} = 1.0$~m, attained at a sample; in this log no interval minimum is below its endpoint values, which is not true in general. (b) With $R = 0.3 + 0.8 = 1.1$~m the pair (drone~2, intruder) collides: $X_{\mathrm{coll}} = 1$, $n_{\mathrm{coll}} = 1$; with an intruder radius of $0.3$~m there is no collision. (c) $L_1 = 3$~m, $L_2 = 1$~m; drone~1 is last away from its goal at $t = 2$, so $T_1 = 3$~s; drone~2 at $t = 1$, so $T_2 = 2$~s; makespan $3$~s, sum of costs $5$~s. (d) $\pos_2 - \pos_1 - \vect{d}_{12}$ is $(0,0), (-1,0), (-1,0), (-2,0)$, so $e_F = 0, 1, 1, 2$~m with mean $1.0$~m and peak $2$~m. (e) The distances are $2, 2.236, 2.236, 2.828$: for $R_{\mathrm{comm}} = 2.5$~m the graph is disconnected ($\lambda_2 = 0$) at $t = 3$ only, one violation; for $2.2$~m at $t = 1, 2, 3$, three violations. \code{\_selftest\_metrics()} in \code{ch25\_evaluation.py} asserts every one of these values.
\end{solution}

\begin{solution}{exr:ch25-critique}
Hints. (1)~Unpaired: ``10 random scenarios per method'' means different scenarios per row; run all methods on the same seeds. (2)~The \orca row is a number copied from a paper: different scenarios, different machine, different implementation; run \orca yourself on the same seeds, or drop the row. (3)~``Averaged over successful runs'' without a success rate hides the failures; add the success rate and the count of unfinished runs, and compare on the commonly finished scenarios. (4)~No spread, no $N$ in the table: add intervals and IQRs. (5)~``$0\,\%$ of 10 runs'' only bounds the rate: the rule of three gives $p \le 0.30$ (exactly $1 - 0.05^{1/10} = 0.26$), the Wilson upper limit is $0.28$; $0\,\%$ and $4\,\%$ are indistinguishable at this $N$. (6)~No minimum separation: two methods with zero collisions may differ greatly in $d_{\min}$. (7)~Absolute path lengths depend on the scenarios; report the ratio to the nominal plan or to the oracle. (8)~Runtimes: state the machine, whether the times are per decision or per run, and add the 99th percentile. (9)~``Ours'' is presumably tuned; state the tuning seeds and whether the baselines were tuned equally. The rewritten table has one row per (cell, method) with $N$, success rate, collision rate [Wilson], $d_{\min}$ median (IQR), $\rho_L$ [interval], makespan, replans, $\bar c$ and $c_{99}$, plus paired-difference rows against the primary baseline.
\end{solution}

\begin{solution}{exr:ch25-pairing}
(a) By \cref{thm:ch25-pairing} with $\sigma_A = \sigma_B = \sigma$, the paired variance is $2\sigma^2(1-\rho)/N_p$ and the unpaired one $2\sigma^2/N_u$; equality needs $N_u = N_p/(1-\rho) = 20/0.1 = 200$ runs per strategy. (b) Pairing hurts when $\rho < 0$: a scenario feature that helps $A$ must hurt $B$, for instance a local avoider that profits from head-on geometry (large closing speed makes the velocity obstacle narrow and early) while a replanner suffers from it. In planner evaluations the dominant scenario features --- density, closing speed, how well aimed the intruder is --- make life harder for every strategy at once, so $\rho$ is positive. (c) Under the null hypothesis each scenario is a fair coin; $\Prob(X \le 4 \mid n = 20) = (1 + 20 + 190 + 1140 + 4845)/2^{20} = 6196/1\,048\,576 = 0.0059$, two-sided $p = 0.0118$, which is the $0.012$ printed for $\rho_L$ of hybrid against replan-only with four drones. (d) In $13$ of $20$ scenarios the hybrid never left its local phase, so its decisions were identical to local-only and the two runs coincide; the seven scenarios in which it differs are exactly those with a replan, and there the hybrid's $d_{\min}$ was larger every time (it never was smaller). The tie count therefore says how often the hybrid's extra machinery was used at all.
\end{solution}

\begin{solution}{exr:ch25-runs}
(a) $(1-p)^N = 0.05$ with $p = 0.005$ gives $N = \ln 0.05 / \ln 0.995 = 597.6$, so $598$ collision-free runs; the rule of three gives $3/0.005 = 600$. A single collision in those runs voids the claim and the count restarts from a larger $N$. (b) The half-width scales with $1/\sqrt{N}$: halving it needs $4 \times 20 = 80$ seeds, and $0.005$ needs $N = 20\,(0.026/0.005)^2 \approx 541$ seeds. (c) The rate is the expensive goal: a mean gains information from every run, while a rate near zero gains it only from the rare events, so a claim about a small rate needs hundreds of clean runs; this is why \cref{sec:ch25-metrics} insists on reporting the minimum separation, which is observed in every run.
\end{solution}
